\noindent\textbf{(a)} En $xy>0$, ponga $u=\sqrt{y/x}>0$ y $\varepsilon=\operatorname{sgn}x$. Las ramas no constantes cumplen $H_\varepsilon(u)=\log|x|+C$, donde\[H_+(u)=-\frac25\log(u+1)-\frac45\log((u-1)^2+1)-\frac45\arctan(u-1),\]\[H_-(u)=-\frac25\log|u-1|-\frac45\log((u+1)^2+1)+\frac45\arctan(u+1).\] Se recupera $y=xu^2$; para $x<0$ se añade $y=x$.

\noindent\textbf{(b)} Con $u=x+y\ne3$,\[y=\frac12\log(u^2-5u+13)+\frac{13}{3\sqrt3}\arctan\frac{2u-5}{3\sqrt3}+C.\]

\noindent\textbf{(c)} $2y+C=\log\left|\dfrac{x-y-1}{x-y+1}\right|$, junto con $y=x-1$ y $y=x+1$. Se exige $x-y\ne0$ y, para la relación logarítmica, $x-y\ne\pm1$.
